\section{Séptima iteración de la Fase 3: Incluir las colisiones en el diseñador gráfico}
\label{sec:ColisionesDisenador}

Las iteraciones anteriores ya \emph{mueven los segmentos}: los del jugador recorren la pantalla mientras los de las rutas avanzan. Falta el corazón del juego: el
momento en que ambos colisionan y la máquina \textbf{decide} que ocurre. Esta iteración declara las colisiones en el \emph{diseñador gráfico}, para que luego el motor de estados pueda consumirlas.

\subsection{Qué se construye}
\label{subsec:QueSeConstruyeF3I7}

\begin{itemize}
	\item \textbf{La colisión como enlace, no como geometría}. La forma clásica de manejar las colisiones es
	calcular el solape de cajas en tiempo real; aquí haremos que el
	autor \textbf{declare} para los segmentos del juego con qué segmento del jugador colisiona (campo
	\texttt{colision}). En ejecución el motor solo comprobará si ese segmento del
	jugador está encendido: una comparación con la bandera que indica \textbf{visible}, sin aritmética de rectángulos (realmente no se requiere siquiera comparar, se puede usar el contenido de la bandera directamente).
	La decisión es del diseñador, no del código: cero ambigüedad de píxeles, y las
	máquinas Game \& Watch reales funcionaban igual (el programador decidía qué
	solapes contaban).
	\item \textbf{El punto de colisión y sus dos desenlaces}. Un segmento que declara
	\texttt{colision} es un \emph{punto de colisión}: una decisión con dos resultados,
	que reutiliza los enlaces que ya existían, sin campos nuevos para la cadena:
	\begin{itemize}
		\item \emph{Colisión positiva}: \texttt{salvado}.
		\item \emph{Colisión negativa}: \texttt{muerte}.
	\end{itemize}
	Cada desenlace sigue el mismo mecanismo usado en este desarrollo: una cadena opcional de segmentos (de juego, de inicio o de vida/muerte) y, ahora, un
	\emph{contador} opcional.
	\item \textbf{Los contadores de desenlace}. Campos nuevos \texttt{salvadoContador} y
	\texttt{muerteContador}: qué contador interno aumenta con la colisión positiva o la
	negativa --- tres valores: \emph{ninguno}, \emph{vida} o \emph{muerte}. Distinción
	importantísima que el análisis dejó cerrada: el \emph{contador interno} de
	vidas/muertes (el que el motor mantiene) no es lo mismo que los segmentos de tipo
	vida (3) o muerte (4), que son su \emph{presentación} en pantalla (los tres
	tiburones bajo el MISS). Y la cadena del desenlace (los enlaces
	\texttt{salvado}/\texttt{muerte}) enlaza con un segmento de juego (la animación
	previa, p. ej. las aletas bajo las barcas, que NO son segmentos de muerte), de
	inicio o de vida/muerte. Sin contador, los contadores internos no se tocan.
	\item \textbf{Reglas por defecto, siempre definidas}. Si un desenlace no lleva
	cadena, el elemento simplemente desaparece; si no lleva contador, los contadores
	quedan intactos. Ninguna colisión puede quedar sin comportamiento: el motor nunca
	tiene que preguntarse qué pasa. El contador interno se toca \textbf{una sola vez},
	en la decisión (no en cada desenlace por separado).
	\item \textbf{Lo que NO es colisión}. Las aletas de los tiburones bajo las barcas en el Parachute, por ejemplo, no
	son puntos de colisión: son el desenlace de una muerte o una animación
	independiente. Eso se resolverá más adelante (segmentos sin
	interacción y secuencias compartidas), y quedan fuera de esta iteración.
	\item \textbf{Un jugador, un punto de colisión}. La regla \emph{un segmento de
		jugador no puede ser punto de colisión de varios elementos} se sostiene: declarar
	una segunda colisión sobre el mismo jugador es un error bloqueante en el propio
	diseñador. Dejarlo múltiple queda como ampliación futura, no como ambigüedad
	presente.
\end{itemize}

\subsection{Detalles de realización}
\label{subsec:DetallesF3I7}

\begin{itemize}
	\item \textbf{Las propiedades nuevas}. \texttt{colision},
	\texttt{salvadoContador} y \texttt{muerteContador} nacen en
	\texttt{propPorDefecto} y entran en \texttt{CAMPOS\_PROP}: viajan en el historial
	(una entrada más del mismo tipo), se deshacen por inversión con su \emph{antes},
	se reproducen al recargar el GWI y --- la colisión, que es un enlace --- se
	re-apunta al jugador resultante cuando una unión absorbe al destinatario,
	exactamente lo que ya hacían los enlaces de cadena y dirección. Cero mecánica
	nueva.
	\item \textbf{El inverso de la colisión}. Igual que la verificación calcula el
	inverso de las direcciones del jugador, ahora se calcula el inverso de la colisión: en
	el segmento del jugador se informa \texttt{colisionDe} (quién choca contra él). Es
	redundancia deliberada, permite simplificar los cálculos para verificar y al generar el GWC, y da al panel del
	jugador su información completa. La unicidad del punto anterior es lo que
	mantiene este inverso unívoco.
	\item \textbf{El panel manda al mapa}. La fila \emph{Colisión con jugador} abre
	la elección en el mapa, como cualquier enlace: clic en el segmento jugador y
	hecho. El diseñador solo acepta el clic si el destino es de tipo jugador, y
	considera la colisión activa en los tipos juego e inicio (los que forman rutas);
	al cambiar de tipo se descarta, como ya ocurría con \texttt{primero}.
	Las filas de los desenlaces viven en el bloque de colisiones, ordenadas como se
	piensan: \emph{Colisión con jugador}, \emph{Colisión positiva (va a)},
	\emph{Colisión negativa (va a)} y los contadores \emph{Contador si positiva} /
	\emph{Contador si negativa}; TODAS dependen de la colisión declarada --- hasta
	que no se elige el jugador con el que choca no hay decisión, y los desenlaces ni
	se activan. En el jugador, la fila \emph{Le choca} informa en vivo de quién
	declara su colisión contra él (sin esperar a la verificación, que es la que fija
	el inverso en el GWI).
	\item \textbf{El mapa habla en verde azulado}. Los puntos de colisión se pintan en color teal (\texttt{\#008080}) en el mapa y en la vista del aparato, con un trazo
	del mismo color entre el punto y su jugador cuando cualquiera de los dos está
	seleccionado: la relación se ve de un vistazo sin leer el panel.
	\item \textbf{Verificación inversa}. El inverso se calcula y se guarda; un
	jugador sin puntos de colisión produce un aviso (la partida sería injugable);
	una colisión cuyo jugador no existe o no es de tipo jugador es \textbf{error
	bloqueante} (no se guarda el GWI); una cadena de desenlace, que enlace con un
	segmento de jugador, también lo es. Si faltara un desenlace, no pasa
	nada grave: rige la	regla por defecto (desaparece sin tocar contadores), y eso queda como aviso, no como error.

\end{itemize}

\begin{figure}[H]
	\centering
	\includegraphics[width=0.7\linewidth]{imagenes/propias/C05-I024}
	\caption[]{Definir colisiones en el diseñador.}
	\label{fig:c05-i024}
\end{figure}
\subsection{Verificación}
\label{subsec:VerificacionF3I7}

\begin{itemize}
	\item La simulación (DOM falso, canvas de píxeles) ejercita el flujo
	completo: declarar una colisión entre un elemento de ruta y el jugador, comprobar
	que el panel solo acepta un destino jugador, que el segundo elemento que apunta al
	mismo jugador se rechaza con error, y que el inverso aparece en el jugador tras
	verificar.
	\item El historial se prueba donde más crítico es: deshacer la declaración de la
	colisión (y su rehacer) recupera el estado exacto, y un GWI guardado con
	colisiones recarga la autoría completa por replay del historial.
	\item Las uniones se prueban con el caso peor: unir dos jugadores (o un jugador
	con otro segmento) re-apunta la colisión al grupo resultante, sin referencias
	colgantes; la verificación posterior lo confirma.
	\item El GWI saliente se inspecciona a mano: \texttt{colision} viaja como
	\texttt{seg\_NNN} y los contadores \texttt{salvadoContador}/\texttt{muerteContador}
	como texto (\texttt{vida}/\texttt{muerte}), solo cuando están puestos, y
	\texttt{colisionDe} en el
	jugador; un GWI sin colisiones no lleva los campos (compatibilidad hacia atrás
	total).
	\item Prueba en navegador real con el cartucho de prueba: el segmento 21
	colisiona con el 23; si es positiva no va a ningún lugar (el elemento
	desaparece, regla por defecto) y aumenta el contador de vidas; si es negativa
	va al 29 (la animación previa, una cadena de juego) y no toca contadores. El
	panel informa en vivo, el mapa pinta el punto en teal con su trazo al jugador y
	la verificación calcula el inverso.
\end{itemize}

\clearpage

